\section*{Scope and verification}

This blueprint follows Yimu Yin's
\href{https://arxiv.org/abs/1006.2467v1}{\emph{Special transformations in
algebraically closed valued fields}}. Its nodes collect the main
definitions and results from Sections 2--10. Each node's title gives the
corresponding paper number. The section status files
give the finer statement-by-statement coverage and hypotheses.

The graph follows mathematical dependencies at the level of the selected
results. The web version links each Lean declaration to its source.
The checked checkpoint is \texttt{a19e25f} on Lean 4.35.0.
\href{https://github.com/daviddhilbert/special-transformations/actions/runs/42813315566}{GitHub Actions run 28}
built the full library and passed 3,543 headline and axiom reports.
The reports allow \texttt{propext}, \texttt{Classical.choice}, and
\texttt{Quot.sound}. A verified theorem retains the hypotheses displayed
in its Lean type.

For the later sections, work in the paper's pure-characteristic-zero
setting over a small parameter base in a sufficiently saturated model.
When a theorem requires a VF/Gamma-generated base, the Gamma parameters
are actual elements of the value-group quotient. The single absorbing RV
element is \(\rv(0)\). Comparisons involving Mathlib's multiplicative
valuation reverse the paper's additive order.

\setcounter{section}{1}
\section{Preliminaries}

\begin{theorem}[2.3. Quantifier elimination]
\label{thm:qe}
\lean{SpecialTransformations.Milestone0.theorem_2_3_quantifier_elimination}
\leanok
The corrected two-sorted theory of algebraically closed valued fields
admits quantifier elimination in the released VF/RV language.
\end{theorem}
\begin{proof}
\leanok
The \href{https://github.com/daviddhilbert/ACVF-QE/tree/939346b0a99899e22a873479364e3cc5d142731c}{pinned ACVF-QE library}
supplies the sorted elimination theorem and its model-transfer
arguments. The \repoFile{LANGUAGE-ERRATA.md}{language erratum}
records the named RV inverse used by the theorem.
\end{proof}

\begin{definition}[2.4. Guarded RV normal form]
\label{def:normal-form}
\lean{SpecialTransformations.Section3.GuardedRVNormalForm.definition_2_4_guarded_rational_normal_form}
\leanok
Definable relations admit the approved disjunction-of-conjunctions
normal form with rational VF factors under their pole guards and
residue terms under their domain guards.
\end{definition}
\begin{proof}
\leanok
\uses{thm:qe}
Translate the eliminated formula to guarded RV literals over the same
parameters. Retain the original zero branches and compile the guards
into well-sorted formulas.
\end{proof}

\begin{theorem}[2.8. C-minimality]
\label{thm:c-minimality}
\lean{SpecialTransformations.Section3.theorem_2_8_c_minimal}
\leanok
In the nontrivially valued ACVF scope, a unary definable VF set is a
finite Boolean combination of balls, including the singleton and
whole-field conventions.
\end{theorem}
\begin{proof}
\leanok
\uses{thm:qe}
Translate sorted definability through VF representatives, use the
valued-language ball theorem, and transport the resulting balls back
to the released language.
\end{proof}

\section{Structural properties}

The \repoFile{Section3/STATUS.md}{Section 3 status} covers the real and
Gamma parameter forms, including the ball and quotient notation used below.

\begin{lemma}[3.1--3.2. Algebraicity descends along definable families]
\label{lem:algebraic-descent}
\lean{SpecialTransformations.Section3.GammaParameters.lemma_3_1_full_gamma_base,SpecialTransformations.Section3.GammaParameters.corollary_3_2_vf_full_base,SpecialTransformations.Section3.GammaParameters.corollary_3_2_rv_full_base}
\leanok
Algebraicity over every index of a nonempty definable family descends
to the original base under the paper's hypotheses. Apply this to VF
fibers of \(\rv\) and RV fibers of the value map.
\end{lemma}
\begin{proof}
\leanok
Use small-parameter saturation to replace the algebraicity witnesses
over the family by finitely many uniform witnesses. Use complete
invariant preimages for Gamma parameters.
\end{proof}

\begin{lemma}[3.3--3.7. Field algebraicity and exchange]
\label{lem:field-algebraicity}
\lean{SpecialTransformations.Section3.lemma_3_3,SpecialTransformations.Section3.corollary_3_4,SpecialTransformations.Section3.corollary_3_5_sorted,SpecialTransformations.Section3.lemma_3_6_vf,SpecialTransformations.Section3.GammaParameters.lemma_3_6_rv,SpecialTransformations.Section3.GammaParameters.corollary_3_7_gamma_base,SpecialTransformations.Section3.GammaParameters.corollary_3_7_vf_gamma_base}
\leanok
Model-theoretic VF algebraicity implies field algebraicity over the
field generated by the VF parameters. This yields the transcendence
degree comparison and finite image of definable RV-to-VF maps. The
corresponding closure operators have exchange in VF and RV. Adjoining
RV or Gamma parameters preserves the stated VF transcendence, and
adjoining Gamma parameters preserves the stated RV transcendence.
\end{lemma}
\begin{proof}
\leanok
\uses{thm:qe,lem:algebraic-descent}
Away from polynomial roots, RV expressions are locally constant in a
VF variable. This rules out isolated algebraic VF points off the field
algebraic closure. The RV argument uses the residue and value
descriptions of the RV sort.
\end{proof}

\begin{lemma}[3.8--3.10. Finite coding and centers]
\label{lem:finite-coding}
\lean{SpecialTransformations.Section3.Average.lemma_3_8,SpecialTransformations.Section3.lemma_3_9,SpecialTransformations.Section3.ClosedBallCodes.lemma_3_10}
\leanok
Finite definable VF sets admit definable RV codes over the original
parameters. Finite definable families of closed balls admit definable
centers.
\end{lemma}
\begin{proof}
\leanok
Use averages in pure characteristic zero and the strict reduction of
equal-RV classes in the coding induction. For ball families, choose
centers independently of the chosen real code for each ball.
\end{proof}

\begin{lemma}[3.11--3.13. Definable points and ball families]
\label{lem:ball-centers}
\lean{SpecialTransformations.Section3.lemma_3_11,SpecialTransformations.Section3.RVFibreConstancy.corollary_3_12,SpecialTransformations.Section3.BallImaginaries.lemma_3_13}
\leanok
The definable-point and RV fiber constancy results give control of
definable subsets of balls. Over a VF/Gamma-generated base, a finite
definable family of balls has a center function over that base.
\end{lemma}
\begin{proof}
\leanok
\uses{thm:c-minimality,lem:finite-coding}
Use the finite unary ball decomposition to bound the exceptional RV
fibers. For a ball family, pass from proper open balls to their closed
hulls, obtain a common finite hitting set, and average within each
ball.
\end{proof}

\begin{corollary}[3.14. Algebraic closure as a model]
\label{cor:acl-model}
\lean{SpecialTransformations.Section3.corollary_3_14,SpecialTransformations.Section3.GammaParameters.corollary_3_14_gamma_base}
\leanok
For the generated bases and nontrivial value groups specified in 3.14,
the algebraic closure forms an elementary ACVF submodel. The named
real-base version retains the original atomic diagram.
\end{corollary}
\begin{proof}
\leanok
\uses{thm:qe,lem:field-algebraicity}
Identify the VF part with the field algebraic closure. Use the RV
representative theorem and the definable witness criterion to obtain
the full two-sorted elementary substructure.
\end{proof}

\begin{lemma}[3.15--3.18. Value fibers and quotient families]
\label{lem:quotient-families}
\lean{SpecialTransformations.Section3.Values.lemma_3_15,SpecialTransformations.Section3.QuotientFamilies.notation_3_16,SpecialTransformations.Section3.lemma_3_17,SpecialTransformations.Section3.AdditiveCosets.lemma_3_18_1,SpecialTransformations.Section3.ClosedCosetChoice.lemma_3_18_2_closed,SpecialTransformations.Section3.CosetProfileChoice.lemma_3_18_2_generated}
\leanok
The value-fiber finiteness and constancy results apply to definable
quotient families. The additive-coset results give constancy away
from finitely many exceptions and definable representatives under
the ball or generation hypotheses of 3.18.
\end{lemma}
\begin{proof}
\leanok
\uses{thm:c-minimality,lem:ball-centers}
Work with the complete real preimages of the quotient relations.
Unary ball geometry controls the exceptions and the checked
center constructions give the required representatives.
\end{proof}

Remark 3.19 is \repoFile{Section3/RETRACTIONS.md}{retracted}.

\section{Dimensions and categories}

The \repoFile{Section4/STATUS.md}{Section 4 status} records the empty-set,
zero-arity and nonemptiness conventions.

\begin{lemma}[4.1--4.5. VF dimension]
\label{lem:vf-dimension}
\lean{SpecialTransformations.Section4.lemma_4_2_actual_gamma,SpecialTransformations.Section4.lemma_4_3,SpecialTransformations.Section4.lemma_4_4_actual_gamma,SpecialTransformations.Section4.corollary_4_5}
\leanok
VF dimension from finite-to-one witnesses agrees with the injective
presentation and the maximal VF transcendence degree. It satisfies
the fiber-maximum formula and the finite-valued image inequality.
\end{lemma}
\begin{proof}
\leanok
\uses{lem:field-algebraicity,lem:finite-coding}
Use field algebraicity for the rank bound, code the finite choices in
RV, and uniformize the presentations over the original base.
\end{proof}

\begin{proposition}[4.6--4.8. Polydiscs and Zariski dimension]
\label{prop:polydiscs}
\lean{SpecialTransformations.Section4.lemma_4_6,SpecialTransformations.Section4.corollary_4_7,SpecialTransformations.Section4.proposition_4_8_closure}
\leanok
Full VF dimension is detected by an open polydisc in a fiber, and an
included RV polydisc gives a dimension bound. The VF dimension of a
definable set agrees with the Krull dimension of its Zariski closure
in the stated scope.
\end{proposition}
\begin{proof}
\leanok
\uses{thm:c-minimality,lem:vf-dimension}
Use full rank to find the polydisc. Compare the rank of the VF
coordinates with transcendence degree on the affine coordinate ring.
\end{proof}

\begin{lemma}[4.9--4.11. RV rank and residue purity]
\label{lem:rv-rank}
\lean{SpecialTransformations.Section4.rvTupleRank_le_rvDim,SpecialTransformations.Section4.ResiduePurity.lemma_4_10_closure,SpecialTransformations.Section4.lemma_4_11}
\leanok
RV dimension is controlled by the RV algebraic-closure rank. In a
residue fiber the normalized rank agrees with Zariski dimension.
A definable value tuple attains the dimension of the original RV set
on its fiber, with infinity permitted for absorbing coordinates.
\end{lemma}
\begin{proof}
\leanok
\uses{lem:field-algebraicity,lem:quotient-families}
Use exchange, the residue-field interpretation and the value-fiber
dimension argument. Keep the defining base through normalization.
\end{proof}

\begin{definition}[4.12--4.20. VF and RV categories and canonical lifts]
\label{def:lift-categories}
\lean{SpecialTransformations.Section4.lemma_4_13,SpecialTransformations.Section4.lemma_4_19,SpecialTransformations.Section4.lemma_4_20}
\leanok
VF objects are classified by VF and RV-fiber dimension bounds. RV
objects carry companion maps, with finite-to-one maps in the
finite-volume categories. Canonical lift takes the RV inverse image
with the companion labels. The checked commuting-map criteria and
dimension invariance control this construction.
\uses{lem:vf-dimension,lem:rv-rank}
\end{definition}

\begin{proposition}[4.21--4.23. Grothendieck constructions]
\label{prop:grothendieck}
\lean{SpecialTransformations.Section4.VFBoundedKGroup,SpecialTransformations.Section4.RVFiniteBoundedGradedKGroup,SpecialTransformations.Section4.lemma_4_23}
\leanok
Disjoint union and product give the graded positive scissor structures.
Group completion and the additive and semiring congruence criteria
are realized for the specified VF and RV categories.
\end{proposition}
\begin{proof}
\leanok
\uses{def:lift-categories}
Present classes by definable isomorphism and scissor relations, then
apply the checked product and completion constructions.
\end{proof}

\section{RV-pullbacks and special bijections}

\begin{definition}[5.1--5.3. Special bijections]
\label{def:special-bijection}
\lean{SpecialTransformations.Section5.lemma_5_2,SpecialTransformations.Section5.SpecialBijection.remark_5_3}
\leanok
A special bijection composes the canonical label adjunction and
centripetal steps. The unary deformation and inverse-fiber results
give the initial geometry for the induction.
\uses{def:lift-categories}
\end{definition}

\begin{theorem}[5.4. Polynomial preparation]
\label{thm:polynomial-preparation}
\lean{SpecialTransformations.Section5.Paper.theorem_5_4,SpecialTransformations.Section5.theorem_5_4_gamma}
\leanok
A special bijection extending the initial chain can prepare the RV
of a VF polynomial on the specified RV polydisc. Finite families
can be prepared while retaining the earlier contractions.
\end{theorem}
\begin{proof}
\leanok
\uses{def:special-bijection,def:normal-form,lem:finite-coding}
Control the roots on vertical fibers with the uniform obstacle
theorem. Center the finite root families and apply the induction
that reduces the remaining root count.
\end{proof}

\begin{theorem}[5.5--5.6. Deformation to a pullback]
\label{thm:deformation}
\lean{SpecialTransformations.Section5.Paper.theorem_5_5,SpecialTransformations.Section5.Paper.corollary_5_6,SpecialTransformations.Section5.Paper.corollary_5_6_gamma}
\leanok
A definable map to an RV tuple can be contracted after a special
bijection in the stated scope. A definable mixed VF/RV set admits a
same-base special bijection to an RV-pullback.
\end{theorem}
\begin{proof}
\leanok
\uses{thm:polynomial-preparation,def:normal-form}
Prepare the finite list of occurring VF polynomials in the defining
relations. Use the resulting RV keys to control membership and the
full graph of the map.
\end{proof}

\begin{corollary}[5.7. Every VF object has an RV lift presentation]
\label{cor:object-surjectivity}
\lean{SpecialTransformations.Section5.Paper.corollary_5_7,SpecialTransformations.Section5.Paper.corollary_5_7_finite,SpecialTransformations.Section5.Paper.corollary_5_7_gamma,SpecialTransformations.Section5.Paper.corollary_5_7_finite_gamma}
\leanok
The canonical lift is surjective on VF isomorphism classes in the
general and finite-volume categories, with the checked named-Gamma
forms.
\end{corollary}
\begin{proof}
\leanok
\uses{thm:deformation,def:lift-categories,lem:vf-dimension}
Construct a same-base mixed-dimensional presentation and deform it
to an RV-pullback. In the finite case, transport finite RV fibers
through the isomorphism.
\end{proof}

\section{Henselian fields}

\begin{definition}[6.1. Functional closure]
\label{def:functional-closure}
\lean{SpecialTransformations.Section6.FunctionallyClosed.over}
\leanok
A substructure is functionally closed when each definable partial
function over it sends tuples from it to tuples in it whenever the
function is defined. This is equivalent to closure under definable
singletons.
\end{definition}

\begin{lemma}[6.2. Henselian substructures are functionally closed]
\label{lem:henselian-closure}
\lean{SpecialTransformations.Section6.Paper.lemma_6_2}
\leanok
Let a substructure of the given pure-characteristic-zero ACVF model
have a nontrivially valued henselian VF field. If every RV element of
the substructure is represented by that field, the substructure is
functionally closed.
\uses{def:functional-closure}
\end{lemma}
\begin{proof}
\leanok
\uses{thm:qe}
Take the algebraic closure of the VF subfield inside the ambient
field and restrict the valuation. Its leading-term model embeds
elementarily. Henselian valuation invariance and the Galois fixed-field
theorem descend definable VF singletons. Averaging the conjugates of
an RV representative descends the RV singletons. See the
\repoFile{Section6/HenselianAlgebraicClosure.lean}{standalone proof}.
\end{proof}

\begin{lemma}[6.3. Elimination of a finite existential VF tuple]
\label{lem:henselian-one-step}
\lean{SpecialTransformations.Section6.Paper.lemma_6_3}
\leanok
In an arbitrary RV-only expansion of the pure-characteristic-zero
henselian theory, eliminate a finite existential VF tuple from a
well-sorted VF-quantifier-free formula. One equivalent formula
preserves the original free VF and RV variables in every model.
\end{lemma}
\begin{proof}
\leanok
\uses{thm:qe,thm:polynomial-preparation}
Use the checked RV literal normalization and uniform vertical
preparation. Construct the valued algebraic envelopes, compare
existential realizations there, and use compactness to choose the
formula uniformly. The \repoFile{Section6/HenselianOneStep.lean}{standalone proof}
supplies these inputs internally.
\end{proof}

\begin{proposition}[6.4. Relative QE for henselian fields]
\label{prop:henselian-qe}
\lean{SpecialTransformations.Section6.Paper.proposition_6_4}
\leanok
The named-base henselian theory admits elimination of VF quantifiers
with arbitrary additional RV-only symbols and the original constants.
\end{proposition}
\begin{proof}
\leanok
\uses{lem:henselian-one-step}
Use the uniform one-step elimination argument recursively and retain
the original names when passing to the atomic diagram.
\end{proof}

\begin{corollary}[6.5. Denef--Pas elimination]
\label{cor:denef-pas}
\lean{SpecialTransformations.Section6.Paper.corollary_6_5}
\leanok
Ordinary pure-characteristic-zero henselian fields with an angular
component admit field-sort elimination in the specified sorted
Denef--Pas expansions. This covers both valuation branches.
\end{corollary}
\begin{proof}
\leanok
\uses{prop:henselian-qe}
Identify the nontrivial branch with the native RV presentation
provided by the angular component. Use the checked direct
translation for the trivial branch.
\end{proof}

\section{Lifting RV correspondences}

\begin{lemma}[7.1--7.2. Weighted polynomials and Hensel lifting]
\label{lem:weighted-hensel}
\lean{SpecialTransformations.Section7.encoded_generalized_hensel}
\leanok
A simple common residue root of a square system of Gamma-polynomials
has a unique simultaneous VF zero in the prescribed RV fiber.
\end{lemma}
\begin{proof}
\leanok
Normalize the weighted polynomials to integral coefficients.
The residue Jacobian is a unit, so multivariate Hensel lifting
gives existence and uniqueness on the normalized fiber.
\end{proof}

\begin{lemma}[7.4--7.5. Polynomial witnesses and constructibility]
\label{lem:residue-witnesses}
\lean{SpecialTransformations.Section7.lemma_7_4,SpecialTransformations.Section7.remark_7_5}
\leanok
Over the stated VF/Gamma-generated base, the RV algebraicity
condition in 7.4 yields a Gamma-polynomial residue-root witness
with nonvanishing last derivative. Definable residue sets are
constructible over the specified residue parameter field.
\end{lemma}
\begin{proof}
\leanok
\uses{lem:rv-rank}
Normalize the residue coordinates over the original base. The
separability argument gives the last-derivative condition, and
residue purity and the released ACF QE give constructibility.
\end{proof}

\begin{theorem}[7.6. Lifting correspondences]
\label{thm:correspondence-lift}
\lean{SpecialTransformations.Section7.theorem_7_6,SpecialTransformations.Section7.theorem_7_6_sorted}
\leanok
In the stated generated-base scope, a definable finite-to-finite
correspondence on nonzero RV tuples has a same-base definable lift.
It restricts to a bijection on each specified pair of RV polydiscs.
\end{theorem}
\begin{proof}
\leanok
\uses{lem:weighted-hensel,lem:residue-witnesses}
Construct local polynomial systems in both directions, use finite
definable covers to select the systems, and apply the unique-root
theorem on each polydisc.
\end{proof}

\begin{corollary}[7.7. Surjectivity on positive Grothendieck semigroups]
\label{cor:unmeasured-scissors}
\lean{SpecialTransformations.Section7.corollary_7_7,SpecialTransformations.Section7.corollary_7_7_finite,SpecialTransformations.Section7.Paper.corollary_7_7_gamma,SpecialTransformations.Section7.Paper.corollary_7_7_finite_gamma}
\leanok
Canonical lift induces surjective homomorphisms on the positive
Grothendieck semigroups in the general and finite-volume cases,
including the checked named-Gamma forms.
\end{corollary}
\begin{proof}
\leanok
\uses{cor:object-surjectivity,thm:correspondence-lift,prop:grothendieck}
The lifted isomorphisms and the behavior on scissor decompositions
give well-defined homomorphisms. Object surjectivity supplies
preimages for generators.
\end{proof}

\section{Local structure of definable maps}

\begin{lemma}[8.1--8.2. Separated sets and isolated fibers]
\label{lem:isolated-fibers}
\lean{SpecialTransformations.Section8.lemma_8_1_gamma,SpecialTransformations.Section8.lemma_8_2_gamma_base}
\leanok
A definable VF set with pairwise disjoint open polydiscs of a common
radius at its distinct points is finite. The isolated points in the
fiber neighborhoods of the stated VF maps form a definable exceptional
locus of dimension below the positive input arity.
\end{lemma}
\begin{proof}
\leanok
\uses{thm:c-minimality,lem:vf-dimension}
Use unary ball geometry and induction for the separated sets. Use
the definable two-radius relation and the dimension comparison for
the isolated-fiber locus.
\end{proof}

\begin{lemma}[8.3. Unary slopes]
\label{lem:unary-slopes}
\lean{SpecialTransformations.Section8.lemma_8_3_canonical,SpecialTransformations.Section8.lemma_8_3_gamma_base}
\leanok
Away from a finite exception in each unary fiber, the local value
behavior of the specified definable map is either constant or given
by an exact value equation with slope in a uniform finite set of
positive rationals.
\end{lemma}
\begin{proof}
\leanok
\uses{def:normal-form,thm:c-minimality}
Choose the candidate slopes from the polynomial value comparisons.
Use unary definability and the checked parameter descent to choose
the local radius and constants.
\end{proof}

\begin{lemma}[8.4--8.5. Slope one and the RV difference law]
\label{lem:rv-difference}
\lean{SpecialTransformations.Section8.lemma_8_4_named_gamma_base,SpecialTransformations.Section8.lemma_8_5_partial_gamma_base}
\leanok
For the centered ball bijections of 8.4, the local slope is one.
The balls have finite radii. For the unary bijections of 8.5,
outside a finite definable exception there is a ball and a nonzero
RV scalar on which the RV difference law holds. Both results take
the explicit henselization premises recorded in
\repoFile{Section8/STATUS.md}{Section 8}.
\uses{hyp:immediate-henselization}
\end{lemma}
\begin{proof}
\leanok
\uses{lem:unary-slopes}
Apply the permitted Engler--Prestel henselization facts over the
field obtained by adjoining a test point. The map and inverse have
integral germ slopes, forcing slope one. Apply that result to the
local difference calculation in 8.5.
\end{proof}

\begin{lemma}[8.6--8.9. Generic estimates, continuity and constancy]
\label{lem:local-regularity}
\lean{SpecialTransformations.Section8.lemma_8_6_gamma_base,SpecialTransformations.Section8.lemma_8_7_gamma_base,SpecialTransformations.Section8.rvLocallyConstantOn,SpecialTransformations.Section8.lemma_8_9_gamma_base}
\leanok
The VF estimates and continuity statements hold off the stated
lower-dimensional loci. The definable RV families of 8.9 are
locally constant on an open VF locus of smaller-dimensional
complement. In 8.6 the base names a positive Gamma element.
\end{lemma}
\begin{proof}
\leanok
\uses{lem:isolated-fibers,lem:unary-slopes,lem:quotient-families}
Uniformize the unary estimates in families and apply the dimension
comparison to the exceptional sets. Use the additive-coset
constancy result in the unary step for RV families.
\end{proof}

\section{Limits, derivatives and Jacobians}

\begin{definition}[9.1. Limit sets]
\label{def:limit-set}
\lean{SpecialTransformations.Section9.IsLimitSet,SpecialTransformations.Section9.IsMinimalLimitSet}
\leanok
For every \(\epsilon\in\Gamma\) there is a \(\delta\in\Gamma\) such
that, for all \(c\in\mathfrak{o}(a,\delta)\cap(A\setminus\{a\})\),
\[
f(c)\subseteq\bigcup_{b\in L}\mathfrak{o}(b,\epsilon).
\]
A limit set is minimal if every proper subset fails to be a limit
set. The \repoFile{Section9/LIMIT-SET-ERRATUM.md}{erratum} records
the equivalent full-cover form and the local boundedness consequence
for finite affine limit sets. The auxiliary point at infinity uses
complements of closed balls.
\end{definition}

\begin{lemma}[9.2--9.5. Finite limit sets]
\label{lem:finite-limits}
\lean{SpecialTransformations.Section9.lemma_9_2,SpecialTransformations.Section9.lemma_9_3,SpecialTransformations.Section9.lemma_9_4,SpecialTransformations.Section9.lemma_9_5}
\leanok
The minimal affine limit sets of 9.2 have the same closure, and
closed minimal limits obey the union rule. A finite minimal limit
set injects into one nearby fiber under the hypotheses of 9.4.
A definable unary relation with finite fibers has a finite
same-base definable minimal limit set in the extended space.
\uses{def:limit-set}
\end{lemma}
\begin{proof}
\leanok
\uses{thm:c-minimality}
Separate finite candidate points by neighborhoods. Obtain the
finite candidates from polynomial roots in the affine and
reciprocal charts, then use the definable frequent-hit relation
for minimality and same-base definability.
\end{proof}

\begin{lemma}[9.6--9.10. Differentiation on VF charts]
\label{lem:vf-differentiation}
\lean{SpecialTransformations.Section9.HasVFDerivativeAt,SpecialTransformations.Section9.HasVFDerivativeAt.comp,SpecialTransformations.Section9.lemma_9_8,SpecialTransformations.Section9.corollary_9_9,SpecialTransformations.Section9.HasMixedVFJacobianAt}
\leanok
The VF derivative is the linear approximation on an open VF chart.
It satisfies the chain rule. The partial derivatives have same-base
definable graphs and are continuous away from a common
lower-dimensional exception in the stated scope. The mixed Jacobian
is the determinant on a chart with RV labels fixed.
\end{lemma}
\begin{proof}
\leanok
\uses{lem:finite-limits,lem:local-regularity}
Use finite limits of difference quotients and differentiation at
independent VF points to obtain the generic results. The chain rule
follows from the linear approximation estimates.
\end{proof}

\begin{lemma}[9.11--9.12. Corrected Jacobian rules]
\label{lem:vf-jacobian-rules}
\lean{SpecialTransformations.Section9.lemma_9_11_ae,SpecialTransformations.Section9.mixedVFJacobian_comp_of_full_derivative,SpecialTransformations.Section9.lemma_9_12_ae}
\leanok
A special bijection and its inverse have determinant one almost
everywhere. On fully differentiable composable fixed-label VF charts,
the Jacobians multiply. The generic composition form has the
finite-to-one first-map hypothesis. The
\repoFile{Section9/DIFFERENTIATION-ERRATA.md}{errata} specify the
pointwise and almost-everywhere scopes.
\end{lemma}
\begin{proof}
\leanok
\uses{def:special-bijection,lem:vf-differentiation}
A centripetal step is a shear on the good charts and canonical
steps have identity VF map there. Apply the chain rule, using
dimension bounds to control the exceptional sets in the composition.
\end{proof}

\begin{definition}[9.13--9.14. Essential maps and RV Jacobians]
\label{def:rv-jacobian}
\lean{SpecialTransformations.Section9.rvEssentialIso,SpecialTransformations.Section9.IsRVJacobianAt,SpecialTransformations.Section9.exists_definable_generic_RVJacobian,SpecialTransformations.Section9.rvJacobian_chain_ae_finite_volume}
\leanok
Essential restriction removes volume tuples with absorbing
coordinates. The generic RV Jacobian is a residue tangent
determinant normalized by the RV scale factors. The
\repoFile{Section9/RV-CHAIN-ERRATUM.md}{corrected chain rule}
has a same-base lower-dimensional exception in positive degree.
\uses{def:lift-categories,lem:rv-rank}
\end{definition}

\begin{proposition}[Section 9 support. Henselized graph-field construction]
\label{prop:henselized-graph-field}
\lean{SpecialTransformations.Section9.acvf_henselizedGaussGraphInput_of_finite_extension_reference}
\leanok
The graph-field data used by the Section 9 Jacobian arguments can be
constructed in the ambient ACVF field under the finite henselian
extension hypothesis.
\uses{hyp:finite-henselian-extension}
\end{proposition}
\begin{proof}
\leanok
\uses{import:stacks-04GN,import:stacks-04GS}
Construct the ordinary henselization of the Gauss valuation ring
and embed it into the ambient valuation ring using the local map
from its universal property. The finite-extension hypothesis supplies
the henselianity and residue-degree calculation for the graph field.
\end{proof}

\begin{lemma}[9.15. Comparing VF lifts with RV Jacobians in positive arity]
\label{lem:rv-jacobian-comparison}
\lean{SpecialTransformations.Section9.lemma_9_15_finite_volume_of_reference_input}
\leanok
For the corrected RV clause, an arbitrary definable lift of a
finite-volume essential RV isomorphism has the specified RV
Jacobian comparison off the indicated lower-dimensional loci.
The positive-arity theorem uses the explicit finite henselian
extension input described below.
\uses{def:rv-jacobian,lem:vf-differentiation,prop:henselized-graph-field}
\end{lemma}
\begin{proof}
\leanok
\uses{thm:correspondence-lift}
Compare the derivatives over the residue-generic graph field,
then descend the result to the original base. The
\repoFile{Section9/LEMMA-9.15.md}{withdrawn Gamma clause}
has a documented valued-field counterexample. Lean checks its
underlying group permutation and change of the sum of valuations.
\end{proof}

\begin{theorem}[9.16. A lift with prescribed Jacobian in positive arity]
\label{thm:prescribed-lift}
\lean{SpecialTransformations.Section9.theorem_9_16_finite_volume_of_reference_input}
\leanok
In the finite-volume generated-base setting, the two compatibility
conditions on the RV isomorphism give a same-base definable VF lift
with the prescribed RV Jacobian away from a VF-small exception.
The proof takes the henselized graph-field data as input.
\uses{def:rv-jacobian,prop:henselized-graph-field}
\end{theorem}
\begin{proof}
\leanok
\uses{thm:correspondence-lift,lem:rv-jacobian-comparison}
Construct local graph and tangent data on the good locus, choose
the compatible definable lift, and assemble the pieces with the
checked exceptional-set bounds.
\end{proof}

\begin{theorem}[9.16. A lift with prescribed Jacobian in arity zero]
\label{thm:prescribed-lift-zero}
\lean{SpecialTransformations.Section9.exists_prescribed_lift_zero_arity}
\leanok
In arity zero the prescribed Jacobian is one and the label
isomorphism lifts in VF.
\uses{def:rv-jacobian}
\end{theorem}
\begin{proof}
\leanok
\uses{thm:correspondence-lift}
Lift the label map using the empty VF coordinate tuple. Its
determinant is the empty product.
\end{proof}

\begin{proposition}[Section 9 support. Gamma descent on RV polydiscs]
\label{prop:gamma-polydisc-descent}
\lean{SpecialTransformations.Section9.gammaDescent_polydiscs_of_vfSmall_exception}
\leanok
For the stated definable bijections of RV polydiscs, a fixed finite
Jacobian value off a VF-small set equals the ratio of the products of
the target and source values in the multiplicative valuation convention.
\uses{prop:henselized-graph-field}
\end{proposition}
\begin{proof}
\leanok
\uses{lem:vf-dimension}
Use the intrinsic coarsening and compare the residue-generic tangent
determinant on the graph field. The dimension bound puts independent
VF tuples outside the exceptional set.
\end{proof}

\section{Categories with volume forms}

\begin{definition}[10.1 and 10.3. Measured categories]
\label{def:volume-categories}
\lean{SpecialTransformations.Section10.vfDegreeCoproductCategory,SpecialTransformations.Section10.rvVolumeGroupoid,SpecialTransformations.Section10.gammaVolumeGroupoid}
\leanok
In the finite-volume setting the VF and RV groupoids carry
definable RV-valued or Gamma-valued forms. VF morphisms are
measured germs. The graded constructions include degree zero
and the positive constant-volume cases.
\uses{def:lift-categories,def:rv-jacobian}
\end{definition}

\begin{remark}[10.2. Almost-everywhere agreement]
\label{rem:measured-germs}
\lean{SpecialTransformations.Section10.MeasuredVF.ae_iff_agreement,SpecialTransformations.Section10.MeasuredVF.quotientEquiv}
\leanok
Agreement on the VF tuples independent over the original
parameter field is equivalent to agreement off a same-base
VF-small exception. The quotient of representatives is
equivalent to the measured germ type.
\uses{def:volume-categories}
\end{remark}
\begin{proof}
\leanok
\uses{lem:vf-dimension,lem:vf-jacobian-rules}
Use the dimension characterization of independent tuples.
Restrict representatives to the good composable overlaps and
apply the Jacobian rules on the actual fixed-label charts.
\end{proof}

\begin{theorem}[10.4. Surjectivity on measured objects]
\label{thm:measured-objects}
\lean{SpecialTransformations.Section10.theorem_10_4_rv_volume,SpecialTransformations.Section10.theorem_10_4_gamma_volume,SpecialTransformations.Section10.theorem_10_4_zero}
\leanok
In the original-base finite-volume scope, a measured VF object has
an essential measured isomorphism to an RV lift. This holds in
both RV-valued and Gamma-valued cases, with the unmeasured
specialization in degree zero.
\uses{def:volume-categories}
\end{theorem}
\begin{proof}
\leanok
\uses{cor:object-surjectivity,thm:deformation,lem:vf-jacobian-rules}
Prepare the form along with the underlying VF set. The constructed
representative has Jacobian one on the good locus, so the prepared
form descends to the RV object.
\end{proof}

\begin{theorem}[10.5. Lifting RV-valued measured isomorphisms]
\label{thm:measured-isomorphisms-rv}
\lean{SpecialTransformations.Section10.theorem_10_5_rv_volume}
\leanok
In positive degree, an essential RV-valued measured isomorphism admits a same-base VF
representative with the given label map and the measure equation
on its actual restricted chart. The source and target complements
have VF dimension below the degree.
\uses{def:volume-categories,rem:measured-germs,hyp:finite-henselian-extension}
\end{theorem}
\begin{proof}
\leanok
\uses{thm:prescribed-lift,lem:vf-jacobian-rules}
Use the lift with prescribed RV Jacobian, then restrict away
from the closure of the exceptional VF projection. The restriction
preserves the determinant on the resulting chart.
\end{proof}

\begin{theorem}[10.5. Lifting Gamma-valued and degree-zero isomorphisms]
\label{thm:measured-isomorphisms-gamma}
\lean{SpecialTransformations.Section10.theorem_10_5_gamma_volume,SpecialTransformations.Section10.theorem_10_5_zero}
\leanok
In positive degree the Gamma-valued form gives a same-base VF
representative with the given label map and the Gamma measure
equation, with source and target complements of dimension below the
degree. Degree zero uses the unmeasured isomorphism.
\uses{def:volume-categories,rem:measured-germs}
\end{theorem}
\begin{proof}
\leanok
\uses{thm:correspondence-lift,lem:vf-jacobian-rules}
Use the checked weighted lift and the pointwise Gamma condition,
then transfer the determinant to the restricted chart. The
degree-zero statement lifts the unmeasured label map.
\end{proof}

\begin{corollary}[10.6. Surjectivity for RV-valued volume]
\label{cor:measured-scissors-rv}
\lean{SpecialTransformations.Section10.corollary_10_6_rv,SpecialTransformations.Section10.corollary_10_6_rv_unit}
\leanok
Canonical measured lift is a surjective additive homomorphism on the
positive scissor presentations for RV-valued volume, including the
constant-volume specialization.
\end{corollary}
\begin{proof}
\leanok
\uses{thm:measured-objects,thm:measured-isomorphisms-rv,prop:grothendieck}
Lift respects measured isomorphism by the RV-valued case of 10.5,
and respects the empty object and cutting relations. Apply 10.4
to each target generator.
\end{proof}

\begin{corollary}[10.6. Surjectivity for Gamma-valued volume and degree zero]
\label{cor:measured-scissors-gamma}
\lean{SpecialTransformations.Section10.corollary_10_6_gamma,SpecialTransformations.Section10.corollary_10_6_gamma_unit,SpecialTransformations.Section10.corollary_10_6_zero}
\leanok
Canonical measured lift is a surjective additive homomorphism on
the positive scissor presentations for Gamma-valued volume and its
constant-volume specialization, and on the unmeasured degree-zero
presentations.
\end{corollary}
\begin{proof}
\leanok
\uses{thm:measured-objects,thm:measured-isomorphisms-gamma,prop:grothendieck}
Use the corresponding isomorphism lifts to pass to scissor classes.
Apply object surjectivity to the target generators.
\end{proof}

The \repoFile{Section10/STATUS.md}{Section 10 status} records the
finite-to-one volume maps, the original generated-base conditions
and the exact degree and constant-volume specializations. The
broader ambient categories in Sections 9--10 lie outside the
approved scope.

\section*{Reference hypotheses and checked imports}

The amber graph nodes mark explicit mathematical hypotheses in the
Lean results. Purple nodes mark the checked Stacks imports. The
ordinary model assumptions, such as characteristic, saturation and
the parameter base, are stated with the individual results and in
the linked section status files.

\begin{assumption}[Immediate henselization]
\label{hyp:immediate-henselization}
Every valued subfield of the ambient field has a henselian extension
inside that field with the same values. Automorphisms of its algebraic
closure over the extension preserve valuation comparisons. This is the
reference premise of Lemmas 8.4 and 8.5.
\repoFile{Section8/STATUS.md}{Section 8 status} cites
\href{https://link.springer.com/book/10.1007/3-540-30035-X}{Engler--Prestel,
\emph{Valued Fields}}, Theorems 4.1.3, 5.2.2 and 5.2.5.
\end{assumption}

\begin{assumption}[Finite henselian extensions]
\label{hyp:finite-henselian-extension}
In residue characteristic zero, a finite extension of a henselian
valued field with unchanged value group is henselian and its residue
degree equals its field degree. This Engler--Prestel premise supplies
the graph-field input for positive-arity Lemma 9.15 (RV clause),
Theorem 9.16 and the intrinsic Gamma polydisc descent. The
positive-degree RV-valued cases of 10.5 and 10.6 also use it.
\repoFile{Section9/GaussGraphConstruction.lean\#L17}{Lean hypothesis}
\end{assumption}

\begin{checkedimport}[Stacks 04GN. Existence of the ordinary henselization]
\label{import:stacks-04GN}
\leanok
A local ring has a henselian local algebra with the same residue
field and a filtered étale presentation.
\href{https://stacks.math.columbia.edu/tag/04GN}{Stacks Lemma 10.155.1}
\end{checkedimport}
\begin{proof}
\leanok
The pinned library builds the vendored proof of
\repoFile{Section9/VendoredStacks/StacksFrontier/RingTheory/Henselian/Henselization/Existence.lean\#L711}{\texttt{IsLocalRing.exists\_henselization}}.
\end{proof}

\begin{checkedimport}[Stacks 04GS. Universal property of the ordinary henselization]
\label{import:stacks-04GS}
\leanok
A local map from a local ring into a henselian local ring extends
uniquely across the ordinary henselization. The Section 9 construction
uses the extension map and its locality.
\href{https://stacks.math.columbia.edu/tag/04GS}{Stacks Lemma 10.155.6}
\end{checkedimport}
\begin{proof}
\leanok
The pinned library builds the vendored
\repoFile{Section9/VendoredStacks/StacksFrontier/RingTheory/Henselian/HenselizationUniversalProperty.lean\#L375}{local extension map and its properties}.
\end{proof}

The \repoFile{Section9/VendoredStacks/README.md}{vendored source record}
gives the AxiomMath Stacks revision and license. Both imports are used
in the construction of the embedded henselization of the Gauss
valuation ring.
The current Section 6 paper exports supply saturation, algebraic
envelopes and ordinary henselian recognition by checked proofs.
